% U3 WS4 -- Type B: solubility and separation. Block 7.
\documentclass{apchem-worksheet}

\apcunit{3}
\apcunittitle{Properties of Substances and Mixtures}
\apcdoctitle{Worksheet 4 \textbullet\ Solubility and Separation}
\apcsubtitle{CED 3.9, 3.10 \textbullet\ every separation exploits one physical difference}

\begin{document}
\wsheader[For solubility questions, cite the structural feature (a polar
bond, an O--H, a long carbon chain) --- not just the word ``polar.'']

\problem[3.10]{Predict which member of each pair is more soluble in water
and give the structural reason:}
\begin{parts}
  \item \ce{CH3CH2OH} or \ce{CH3CH2CH2CH2CH2OH}
        \answerlines[2]{Ethanol --- both have one O--H, but the longer
        pentanol chain adds nonpolar bulk that water cannot solvate, so the
        polar fraction of the molecule is smaller.}
  \item \ce{I2} or \ce{NaI}
        \answerlines[2]{\ce{NaI} --- its ions are stabilized by strong
        ion--dipole attractions with water; \ce{I2} is nonpolar and offers
        only dispersion forces.}
  \item \ce{CH4} or \ce{NH3}
        \answerlines[2]{\ce{NH3} --- it hydrogen bonds with water;
        \ce{CH4} is nonpolar with only dispersion forces.}
\end{parts}

\problem[3.10]{Vitamin A is composed almost entirely of carbon and hydrogen;
vitamin C contains several O--H and C--O bonds.}
\begin{parts}
  \item Which is water-soluble? Justify from structure.
        \answerlines[2]{Vitamin C --- its polar O--H groups hydrogen bond
        with water (hydrophilic); vitamin A is essentially nonpolar
        (hydrophobic).}
  \item Explain why excess vitamin A can accumulate in the body while
        excess vitamin C does not. \answerlines[2]{Vitamin A dissolves in
        fatty tissue and is stored there; vitamin C stays in the aqueous
        phase and is excreted in urine.}
\end{parts}

\problem[3.10]{Iodine is far more soluble in \ce{CCl4} than in water, even
though water is the more common solvent. Explain using ``like dissolves
like.'' \answerlines[3]{\ce{I2} and \ce{CCl4} are both nonpolar, so the
dispersion forces between them are comparable to those being replaced.
Dissolving \ce{I2} in water would require breaking strong hydrogen bonds
between water molecules and replacing them with much weaker \ce{I2}--water
attractions --- not favorable.}}

\problem[3.9]{Name the technique and the property it exploits:}
\begin{parts}
  \item Separating sand from water: \answerblank[6cm]{filtration ---
        particle size}
  \item Separating ethanol from water: \answerblank[6cm]{distillation ---
        boiling point}
  \item Separating leaf pigments: \answerblank[7.7cm]{chromatography ---
        relative attraction}
  \item Recovering dissolved salt from seawater:
        \answerblank[7.3cm]{evaporation/distillation --- volatility}
\end{parts}

\problem[3.9]{A paper chromatogram is run with a polar stationary phase
(paper) and a nonpolar mobile phase. Two dyes are applied; dye X travels
\SI{7.2}{\centi\meter} and dye Y travels \SI{2.1}{\centi\meter} while the
solvent front travels \SI{9.0}{\centi\meter}.}
\begin{parts}
  \item Which dye is more polar? Justify.
        \answerlines[2]{Dye Y --- it moved less, meaning it was held more
        strongly by the polar paper (stationary phase).}
  \item Compute both retention factors ($R_f = $ distance moved by
        substance $\div$ distance moved by solvent).
        \answerblank[5.5cm]{X: 0.80\quad Y: 0.23}
  \item Predict what happens to dye Y's $R_f$ if a polar mobile phase is
        used instead. \answerlines[2]{It increases --- a polar solvent
        competes for the polar dye and carries it farther up the paper.}
\end{parts}

\problem[3.9]{A mixture contains sand, salt, and water. Design a full
separation, giving the order and the reason for each step.
\answerlines[3]{1.~Filter --- removes sand by particle size. 2.~Distill or
evaporate the filtrate --- separates water from salt by boiling point;
collect the condensed water if it is wanted, leaving solid salt behind.}}

\begin{spiralstrip}{Unit 3 \textbullet\ blocks 1--6}
  \item $[\ce{Cl-}]$ in \SI{0.30}{\Molar} \ce{AlCl3}:
        \answerblank[2.6cm]{\SI{0.90}{\Molar}}
  \item Dilute \SI{20.0}{\milli\liter} of \SI{3.0}{\Molar} to
        \SI{120}{\milli\liter}: \answerblank[2.6cm]{\SI{0.50}{\Molar}}
  \item KMT: average KE depends only on:
        \answerblank[3cm]{temperature}
  \item Ionic solids conduct when: \answerblank[4cm]{molten or dissolved}
  \item Higher bp: \ce{H2O} or \ce{H2S}? Why?
        \answerblank[5.4cm]{\ce{H2O} --- hydrogen bonding}
\end{spiralstrip}

\end{document}
